\documentclass[11pt,a4paper]{article}
\usepackage{../shared/hanzocover}
\usepackage[utf8]{inputenc}
\usepackage[T1]{fontenc}
% Shared preamble for the compute-market series.
% One style, one vocabulary, inputted by every paper so the series reads as one voice.
\usepackage[margin=1in]{geometry}
\usepackage{booktabs}
\usepackage{amsmath}
\usepackage{amssymb}
\usepackage{amsthm}
\usepackage{graphicx}
\usepackage{xcolor}
\usepackage{hyperref}
\hypersetup{colorlinks=true, linkcolor=blue, urlcolor=blue, citecolor=blue}
\usepackage{enumitem}
\usepackage{listings}
\usepackage{array}
\usepackage{tabularx}
\usepackage{siunitx}
\sisetup{group-separator={,}, group-minimum-digits=4}

\definecolor{kw}{rgb}{0.13,0.29,0.53}
\definecolor{cmt}{rgb}{0.40,0.40,0.40}
\definecolor{str}{rgb}{0.16,0.50,0.16}
\lstset{
  basicstyle=\ttfamily\footnotesize,
  keywordstyle=\color{kw}\bfseries,
  commentstyle=\color{cmt}\itshape,
  stringstyle=\color{str},
  showstringspaces=false,
  breaklines=true,
  columns=fullflexible,
  frame=single,
  framesep=4pt,
}

\newtheorem{definition}{Definition}
\newtheorem{proposition}{Proposition}
\newtheorem{remark}{Remark}

\newcommand{\code}[1]{\texttt{#1}}
\newcommand{\repo}[1]{\texttt{#1}}

% Series vocabulary — used identically in every paper.
\newcommand{\job}{\ensuremath{J}}
\newcommand{\bundle}{\ensuremath{R}}
\newcommand{\cost}{\ensuremath{\mathcal{C}}}

\tolerance=800
\emergencystretch=3em


\title{The Hamiltonian Market Maker:\\Regime-Aware Pricing for Perishable Heterogeneous Compute}

\author{
Hanzo AI Research\thanks{Corresponding author: research@hanzo.ai} \\
\textit{Hanzo Industries Inc (Techstars '17)} \\
Los Angeles, California \\
\texttt{https://hanzo.ai}
}

\date{2026}

\begin{document}
\hanzocoverpage

\begin{abstract}
Constant-product market makers price a fungible, storable pair. Compute is
none of those things: a GPU-second not sold at 14:00 is gone, an H100-second and
a CPU-second are not substitutes at any exchange rate, and a job with a
three-second deadline values capacity differently from the same job with an hour.
An $x \cdot y = k$ curve has no way to express scarcity convexity, deadline
pressure, or per-resource isolation, and a market that adopts one inherits those
silences.

We describe an alternative that we have implemented and run: price as position
and trend as momentum in a phase space, evolved by an energy-conserving
integrator under a regime-dependent anharmonic potential, with the regime itself
recovered as the latent state of a hidden Markov model driven by request rate,
queue depth and utilisation. Each resource kind carries its own market state and
prices never cross between kinds, which is the property that makes heterogeneity
representable rather than averaged away. Per-tenant preference is carried as a
one-bit sign delta over a base quote table, small enough that millions of tenant
tables sit in memory at once.

The result is a quote that is deterministic given a job and a market state ---
reproducible and testable rather than emergent --- and that responds to a demand
spike by changing its own dynamics rather than by sliding along a fixed curve.
\end{abstract}

\tableofcontents
\newpage

\section{What a curve cannot say}

A constant-product maker holds reserves $(x, y)$ and quotes the price implied by
$x \cdot y = k$. The design is superb for what it was built for, and each of its
assumptions fails for compute.

\paragraph{Fungibility.} The curve assumes any unit of $x$ substitutes for any
other. Accelerator-seconds do not: memory capacity, memory bandwidth, interconnect
and resident state make two nominally identical devices different suppliers for a
given job. Averaging them into one reserve destroys exactly the information the
allocator needs.

\paragraph{Storability.} Reserves persist. Capacity does not. An idle
accelerator-second is not inventory carried to the next block; it is destroyed.
Perishability means the reservation price falls as the window closes, which no
static curve represents.

\paragraph{Deadline neutrality.} The curve prices quantity, not urgency. Two
requests for identical resources with different service objectives are different
products with different willingness to pay, and the difference is not a quantity
difference.

\paragraph{Scarcity convexity.} Near exhaustion, the marginal value of the last
unit of a specific resource rises faster than any fixed hyperbola. A curve chosen
in advance either underprices the tail or overprices the middle, permanently.

\section{Phase space}

We take a different primitive. Let $q$ denote demand pressure and $p$ the shadow
price of a resource kind, and treat the pair as a point in phase space evolving
under a Hamiltonian
\begin{equation}
H(q,p) \;=\; \frac{p^{2}}{2m} \;+\; V_{\varrho}(q),
\end{equation}
where $V_{\varrho}$ is a potential whose shape depends on the market regime
$\varrho$. The equations of motion are the usual ones, with a dissipative term so
the system settles rather than ringing forever:
\begin{equation}
\dot q \;=\; \frac{\partial H}{\partial p} \;=\; \frac{p}{m},
\qquad
\dot p \;=\; -\frac{\partial H}{\partial q} \;-\; \gamma\, p .
\end{equation}

Price is position; trend is momentum. That is the whole of the intuition, and it
buys three things a curve does not have.

\paragraph{Momentum is a state variable, not an inference.} A market that has
been rising carries that fact in $p$. A curve has to rediscover trend from
reserve history every time it is asked.

\paragraph{Integration is energy-conserving.} We integrate symplectically. The
practical consequence is that the maker does not drift on noise: a symplectic
integrator has no secular energy error, so a long run of small random shocks does
not accumulate into a phantom trend. A naive Euler step on the same dynamics
gains energy monotonically and would manufacture exactly the drift we are trying
to avoid.

\paragraph{Convexity is in the potential, where it can be changed.} An anharmonic
potential gives the steepening response near exhaustion that a hyperbola cannot,
and because the potential is a separate object it can be re-specified per regime
without touching the integrator.

\section{Regime as a latent variable}

The potential's dependence on $\varrho$ is the second half of the design, and it
is where the other reading of the acronym lives: the regime is the hidden state
of a hidden Markov model.

\begin{definition}[Market regime]
A discrete latent state $\varrho_{t}$ evolving as a Markov chain with learned
transition matrix $A$, emitting observations $o_{t}$ --- request arrival rate,
queue depth, accelerator utilisation. The posterior $\Pr[\varrho_{t} \mid
o_{1:t}]$ is maintained by the forward recursion; the maximum-likelihood
trajectory is available by Viterbi decoding and the transition and emission
parameters are fit by Baum--Welch.
\end{definition}

The states we carry distinguish an exploring market with high posterior
uncertainty, an exploiting market with low uncertainty, a market in crisis under
extreme conditions, and a market in transition between them. The regime modulates
two things the maker controls directly: the width of its quotes and its inventory
bias.

This composition is the substantive claim of the paper. A market maker that
detects its own regime and changes its dynamics accordingly is doing something
categorically different from one that slides along a fixed curve --- it is not
merely repricing, it is reparameterising. A demand spike does not move the quote
along a hyperbola; it changes which potential the quote is evolving under.

\begin{remark}
The two readings of the acronym are not a pun we are stuck with. The hidden
Markov layer is a standalone sequence model, usable for state detection, sequence
prediction and anomaly detection with no market attached; the market maker is a
consumer of it. Keeping them separable is what allows the regime detector to be
tested against synthetic sequences with known ground truth, which is not possible
for a component fused into a pricing loop.
\end{remark}

\section{Heterogeneity by construction}

Each resource kind --- accelerator, processor, memory, storage, bandwidth ---
carries its own market state, and prices never cross between kinds. This is a
property we enforce rather than an emergent outcome, and it is worth being
explicit about why.

A single blended compute price is a lie of composition. It says that a shortage
of memory bandwidth and a shortage of accelerator throughput are the same event,
and it produces quotes that are wrong in both directions simultaneously: too high
for the resource in surplus, too low for the one that is binding. Per-kind
isolation means a bandwidth spike raises the bandwidth quote and leaves the
processor quote alone, which is both correct and legible to a provider trying to
decide what to add.

\begin{table}[h]
\centering
\begin{tabular}{lll}
\toprule
Layer & Object & Property enforced \\
\midrule
Regime detector & hidden Markov model & latent state from observable load \\
Dynamics & symplectic integrator & no drift on noise \\
Potential & anharmonic, per regime & scarcity convexity \\
Modulation & service-objective tightness & perishability, deadline pressure \\
Isolation & per-kind market state & prices never cross kinds \\
Bounds & floor and ceiling & a quote is never free or unbounded \\
\bottomrule
\end{tabular}
\caption{The pricing stack and the property each layer is responsible for.}
\end{table}

\section{Determinism as a design requirement}

The quoting entry point is deterministic for a given job and market state. It
uses the noise-free integrator path exclusively --- no random draw enters a
quote.

This is not fastidiousness. A market whose quotes are irreproducible cannot be
audited, cannot be regression-tested, and cannot settle a dispute about what it
offered. Determinism at the quoting layer means a disputed quote can be
recomputed from the recorded job and market state by anyone, which converts a
disagreement about pricing into an arithmetic check. It also means the equilibrium
the dynamics reach is a testable object rather than a property one observes and
hopes persists.

Where randomness genuinely belongs --- exploration in the adaptive layer --- it
sits outside the quote path, in the routing policy, where its effects are
measurable and can be turned off.

\section{Per-tenant preference at one bit}

Tenants differ in how they trade price against latency and quality, and a maker
that quotes every tenant identically is leaving surplus on the table. The
straightforward implementation --- a full-precision preference table per tenant
--- does not survive contact with scale.

We carry the tenant's deviation as a one-bit sign delta over a shared base quote
table: one bit per regime-and-resource cell, with per-layer scale factors and
higher-precision bias terms retained. The compression against a
single-precision table is thirty-two-fold, which is the difference between
holding thousands of tenant tables in memory and holding millions.

There is a pleasing symmetry here worth naming, because it is the same idea
appearing at two layers of the stack. At the model layer, personalization is a
low-rank or low-bit delta over frozen shared weights. At the pricing layer,
personalization is a one-bit delta over a shared quote table. In both cases the
shared object is large and immovable, the per-tenant object is tiny, and the
economics of the whole system follow from that ratio. The companion paper on
residency develops the model-layer case.

\section{Active inference in the routing decision}

Given a quote, something still has to choose a venue. We score candidate actions
by expected free energy, decomposed into an epistemic term measuring information
gain and a pragmatic term measuring progress toward the requested outcome.

The reason to prefer this over pure expected-cost minimisation is that a market
maker is always operating on an estimate of its own state, and the value of an
allocation includes what it teaches. Routing a job to a provider whose
performance is uncertain has option value that a cost-minimising allocator scores
at zero and then never resolves --- the classic failure where a system converges
on a locally good supplier and never discovers a better one because it never
tries. Making the epistemic term explicit puts exploration in the objective
rather than bolting it on as a heuristic.

The dispatcher's output is a quote --- a price and a venue --- not a model
identity. That separation matters: which model answers is a question for the
layer above, and conflating the two produces a router that cannot be reused when
the model catalogue changes.

\section{Two layers, and why they compose}

A complete market has two allocation problems that are easy to confuse.

The lower problem is \emph{which resources, at what price}. That is this paper:
heterogeneous, perishable, deadline-bound capacity priced by regime-aware
dynamics.

The upper problem is \emph{which intelligence, for this request}. That is a
different question with a different feature space --- request content, task type,
context length, quality requirement --- and a different answer shape. We solve it
with a learned per-request router that featurises the request, scores every
candidate arm by a bilinear utility, and returns the best arm subject to hard
service-objective ceilings, adapting online per tenant. Its fan-out variant
selects the top several arms, dispatches them concurrently, and fuses the
candidate answers with a synthesizer carrying no weights of its own, so that
adding a new arm is adding a row rather than retraining a policy.

The layers compose cleanly because their interfaces are narrow. The upper layer
asks for capacity meeting a specification and does not care how the price was
formed; the lower layer supplies a quote and does not care which model will run.
A market can adopt either without the other, which is the property that makes
them separately useful and jointly more than the sum.

\section{Availability and status}

The pricing crate is roughly forty-seven hundred lines of Rust implementing the
hidden Markov core --- forward--backward, Viterbi, Baum--Welch --- alongside the
Hamiltonian dynamics, the compute-pricing adapter, regime detection, one-bit
tenant adapters, expected-free-energy routing and persistence. It is deliberately
not a model-serving engine; the serving path is a separate component and the
separation is enforced at the crate boundary.

We state two limitations plainly. The parameters --- friction, energy scale,
potential shape per regime, the bounds --- are chosen and defensible rather than
fit against a large live market, because the market that would fit them is what
this work is proposing to build. And the regime taxonomy is coarse; four states
is a starting point, and a market with real load data should expect to refit both
the state count and the transitions rather than inherit ours.

Neither limitation is load-bearing for the argument. The claim is not that our
parameters are right; it is that price-as-position under a regime-dependent
potential is the right \emph{shape} for perishable heterogeneous capacity, that
the shape is implementable, and that we have implemented it.

\end{document}
